\documentclass[11pt,a4paper]{article}
\usepackage[margin=0.7in, top=0.6in, bottom=0.6in]{geometry}
\usepackage[utf8]{inputenc}
\usepackage[T1]{fontenc}
\usepackage{charter}
\usepackage{fontawesome5}
\usepackage{xcolor}
\usepackage{enumitem}
\usepackage{hyperref}
\usepackage{microtype}

\definecolor{accent}{HTML}{1E3A5F}     % Professional Dark Navy
\definecolor{accentblue}{HTML}{2B6CB0} % Slate Blue
\definecolor{lightgray}{HTML}{F3F4F6}
\definecolor{darktext}{HTML}{1F2937}
\definecolor{bulletline}{HTML}{9CA3AF}
\definecolor{subtletext}{HTML}{4B5563}

\hypersetup{colorlinks=true, linkcolor=accentblue, urlcolor=accentblue}

\setlength{\parindent}{0pt}
\setlength{\parskip}{8pt}

\begin{document}
\thispagestyle{empty}
\color{darktext}

% Encabezado idéntico al CV
\begin{center}
    {\Huge \bfseries\color{accent} SANTIAGO ALEJANDRO OCHOA QUESADA} \\[6pt]
    {\small\color{darktext}
    \faEnvelope\ \href{mailto:santy.alejo4@gmail.com}{santy.alejo4@gmail.com}
    \quad\textcolor{bulletline}{\textbar}\quad
    \faPhone\ +57 320 351 7153
    \quad\textcolor{bulletline}{\textbar}\quad
    \faMapMarker*\ Bogotá, Colombia \\[4pt]
    \faLinkedin\ \href{https://linkedin.com/in/santiago-alejandro-ochoa-quesada-a07a18288/}{LinkedIn}
    \quad\textcolor{bulletline}{\textbar}\quad
    \faGithub\ \href{https://github.com/Sochoaq}{GitHub}
    }
\end{center}
\vspace{2pt}
{\color{bulletline}\hrule height 0.8pt}
\vspace{10pt}

% Metadatos de la carta
\noindent
\begin{minipage}[t]{0.62\textwidth}
    \textbf{\color{accent} Para:}\\
    Equipo de Selección \& Talento Tech\\
    \textbf{Experis Colombia} (ManpowerGroup)\\
    Bogotá, Colombia
\end{minipage}
\hfill
\begin{minipage}[t]{0.35\textwidth}\raggedleft
    \textbf{\color{accent} Fecha:}\\
    23 de octubre de 2026
\end{minipage}

\vspace{6pt}
\noindent
{\large\bfseries\color{accent} Asunto: Postulación al cargo de Ingeniero de Inteligencia Artificial (Databricks / GenAI)}
\vspace{2pt}

\noindent Estimado Equipo de Selección de Experis Colombia,

Me dirijo a ustedes con gran entusiasmo para presentar mi postulación a la posición de \textbf{Ingeniero de Inteligencia Artificial}. Como Ingeniero Electrónico de la \textbf{Universidad Nacional de Colombia} (promedio de 4.6/5.0) y actual \textbf{Ingeniero de Proyectos} en AM4, he enfocado mi trayectoria en la intersección entre analítica computacional avanzada, ingeniería de software en Python y construcción de soluciones inteligentes orientadas a resolver retos operativos de negocio.

A lo largo de mi experiencia técnica y proyectos aplicados en entornos industriales y analíticos, he desarrollado competencias clave directamente alineadas con los objetivos de Experis:

\begin{itemize}[label={\color{accentblue}\raisebox{0.2ex}{\tiny\faCircle}}, leftmargin=1.4em, itemsep=3pt, parsep=0pt, topsep=2pt]
    \item \textbf{IA Generativa y Arquitecturas RAG:} Diseño e implementación de sistemas de indexación semántica y recuperación contextual (\textit{Retrieval-Augmented Generation}), así como agentes autónomos con capacidad de ejecución de herramientas (\textit{tool calling}) y orquestación con frameworks como LangChain.
    \item \textbf{Ingeniería de Datos y Databricks:} Desarrollo de pipelines en Python para la ingestión, limpieza, transformación y modelado de datos a gran escala, optimizando el acceso y la trazabilidad de la información operativa.
    \item \textbf{MLOps y Arquitectura de APIs:} Construcción de microservicios robustos con FastAPI, contenedorización en Docker, flujos CI/CD con GitHub Actions y esquemas de evaluación continua para garantizar observabilidad, gobernanza y confiabilidad en producción.
\end{itemize}

Cuento con un nivel de \textbf{inglés B2} (comunicación técnica fluida tanto escrita como conversacional), residencia en Bogotá con plena disponibilidad para esquemas de trabajo híbridos o flexibles, y un enfoque riguroso hacia la gestión de entregables, el cumplimiento de cronogramas y la interacción efectiva con stakeholders multidisciplinarios.

Agradezco de antemano el tiempo dedicado a revisar mi postulación. Me encantaría tener la oportunidad de conversar en una entrevista sobre cómo mi experiencia práctica y pasión por la IA Generativa pueden generar valor inmediato en los proyectos de Experis y ManpowerGroup.

\vspace{10pt}
\noindent Cordialmente,

\vspace{12pt}
\noindent
\textbf{\large\color{accent} Santiago Alejandro Ochoa Quesada}\\
{\small\color{subtletext} Ingeniero Electrónico | AI Engineering \& Software}\\
{\footnotesize\color{subtletext} santy.alejo4@gmail.com \;\textbar\; +57 320 351 7153 \;\textbar\; Bogotá, Colombia}

\end{document}
